% Book 3, Year 1, chapter 24 -- graphs and traversals.
% Provenance (never printed): sources/book3/BRIEFS.md, chapter 24.
% ledger: b1:robbins-1939, b1:kosaraju
\chapter{Graphs and Traversals}\label{ch:b1:graph-traversals}

Brenmoor has thirty places, fifty-three roads and a single one-way street, from the Town Square to the Library. Can
every place still reach every other? The town council would like to go further and make every road one-way, to
widen the pavements: is that possible without cutting anyone off? Both questions are about reachability in a \omterm{def:b1:graph-traversals:graph}{graph},
and both are answered by a program that walks the \omterm{def:b1:graph-traversals:graph}{graph} in a well-chosen order. This chapter builds the two walks,
depth first and breadth first, and what they compute: \omterm{def:b1:graph-traversals:distance}{distances}, components, orders, and the roads that cannot be
spared.

\begin{recall}
The second book drew \omterm{def:b1:graph-traversals:graph}{graphs}, computed \omterm{def:b1:graph-traversals:distance}{distances} in a \omterm{def:b1:graph-traversals:graph}{graph} by \omterm{def:b1:graph-traversals:bfs}{breadth-first search} and routes by Dijkstra's
algorithm. The mathematics book represents a \omterm{def:b1:graph-traversals:graph}{graph} by its adjacency matrix. Stacks and queues were built in
\cref{ch:b1:linked-structures}.
\end{recall}

\section{Definitions}\label{sec:b1:graph-traversals:definitions}

\begin{definition}[Graph, vertex, edge, arc, loop]\label{def:b1:graph-traversals:graph}
A \emph{graph}\index{graph} $G = (V, E)$ is a finite set $V$ of \emph{vertices}\index{vertex} and a set $E$ of
\emph{edges}\index{edge (graph)}. In an \emph{undirected graph} an edge is a set $\{u, v\}$ of two vertices, its
\emph{ends}, written $uv$. In a \emph{directed graph}\index{directed graph} an edge is a pair $(u, v)$, an
\emph{arc}\index{arc} from $u$ to $v$, written $u \to v$. An edge whose two ends are equal is a
\emph{loop}\index{loop (graph)}; unless stated otherwise, graphs here have no loops and $n = |V|$, $m = |E|$.
\end{definition}

\begin{definition}[Degree]\label{def:b1:graph-traversals:degree}
The \emph{degree}\index{degree} $\deg(v)$ of a \omterm{def:b1:graph-traversals:graph}{vertex} of an undirected \omterm{def:b1:graph-traversals:graph}{graph} is the number of \omterm{def:b1:graph-traversals:graph}{edges} at $v$; its
\emph{neighbours} are the other ends of these \omterm{def:b1:graph-traversals:graph}{edges}. In a \omterm{def:b1:graph-traversals:graph}{directed graph}, the \emph{out-degree}\index{out-degree} of
$v$ counts the \omterm{def:b1:graph-traversals:graph}{arcs} leaving $v$, its \emph{in-degree}\index{in-degree} those arriving at $v$.
\end{definition}

\begin{proposition}[Handshake lemma]\label{prop:b1:graph-traversals:handshake}
In an undirected \omterm{def:b1:graph-traversals:graph}{graph}, $\sum_v \deg(v) = 2m$; so the number of \omterm{def:b1:graph-traversals:graph}{vertices} of odd \omterm{def:b1:graph-traversals:degree}{degree} is even. In a \omterm{def:b1:graph-traversals:graph}{directed graph},
the \omterm{def:b1:graph-traversals:degree}{out-degrees} and the \omterm{def:b1:graph-traversals:degree}{in-degrees} both sum to $m$.
\end{proposition}

\begin{proof}
Each edge $uv$ is counted twice in the sum, once in $\deg(u)$ and once in $\deg(v)$; each \omterm{def:b1:graph-traversals:graph}{arc} once among the
\omterm{def:b1:graph-traversals:degree}{out-degrees}, at its tail, and once among the \omterm{def:b1:graph-traversals:degree}{in-degrees}, at its head. The sum of the even \omterm{def:b1:graph-traversals:degree}{degrees} is even, so the
odd \omterm{def:b1:graph-traversals:degree}{degrees} have an even sum, which needs an even number of them.
\end{proof}

\begin{definition}[Path, simple path, cycle]\label{def:b1:graph-traversals:path}
A \emph{path}\index{path (graph)} of length $k$ from $u$ to $v$ is a sequence of \omterm{def:b1:graph-traversals:graph}{vertices} $u = v_0, v_1, \dots, v_k = v$
where each $v_{i-1} v_i$ is an edge (an \omterm{def:b1:graph-traversals:graph}{arc} $v_{i-1} \to v_i$ in a \omterm{def:b1:graph-traversals:graph}{directed graph}). It is \emph{simple}\index{simple
path} if its \omterm{def:b1:graph-traversals:graph}{vertices} are distinct. A \emph{cycle}\index{cycle (graph)} is a path of length at least 1 from a \omterm{def:b1:graph-traversals:graph}{vertex}
to itself whose other \omterm{def:b1:graph-traversals:graph}{vertices} are distinct and distinct from it; in an undirected \omterm{def:b1:graph-traversals:graph}{graph} its length must be at least 3
and it uses no edge twice. $v$ is \emph{reachable} from $u$ if there is a path from $u$ to $v$.
\end{definition}

\begin{definition}[Distance]\label{def:b1:graph-traversals:distance}
The \emph{distance}\index{distance (graph)} $\delta(u, v)$ is the smallest length of a \omterm{def:b1:graph-traversals:path}{path} from $u$ to $v$, counted in
\omterm{def:b1:graph-traversals:graph}{edges}, and $\infty$ if $v$ is not reachable from $u$.
\end{definition}

\begin{definition}[Connected, strongly connected]\label{def:b1:graph-traversals:connected}
An undirected \omterm{def:b1:graph-traversals:graph}{graph} is \emph{connected}\index{connected graph} if every \omterm{def:b1:graph-traversals:graph}{vertex} is reachable from every other. Being
reachable from one another is an equivalence relation on the \omterm{def:b1:graph-traversals:graph}{vertices}; its classes are the \emph{connected
components}\index{connected component (graph)}. In a \omterm{def:b1:graph-traversals:graph}{directed graph}, $u$ and $v$ are \emph{strongly connected} when
each is reachable from the other; the classes of this equivalence are the \emph{strongly connected
components}\index{strongly connected component}, and the \omterm{def:b1:graph-traversals:graph}{graph} is \emph{strongly connected}\index{strongly connected
graph} if there is only one.
\end{definition}

\begin{definition}[Directed acyclic graph]\label{def:b1:graph-traversals:dag}
A \emph{directed acyclic graph}\index{directed acyclic graph}, or DAG, is a \omterm{def:b1:graph-traversals:graph}{directed graph} without a \omterm{def:b1:graph-traversals:path}{cycle}.
\end{definition}

\begin{definition}[Free tree, forest]\label{def:b1:graph-traversals:free-tree}
A \emph{free tree}\index{free tree} is a \omterm{def:b1:graph-traversals:connected}{connected} undirected \omterm{def:b1:graph-traversals:graph}{graph} without a \omterm{def:b1:graph-traversals:path}{cycle}; a \emph{forest}\index{forest} is
an undirected \omterm{def:b1:graph-traversals:graph}{graph} without a \omterm{def:b1:graph-traversals:path}{cycle}, whose components are free trees. Choosing a \omterm{def:b1:graph-traversals:graph}{vertex} as the root turns a free tree
into a \omterm{def:b1:induction-and-termination:rooted-tree}{rooted tree} (\cref{def:b1:induction-and-termination:rooted-tree}).
\end{definition}

\begin{proposition}[Free trees]\label{prop:b1:graph-traversals:free-tree}
For an undirected \omterm{def:b1:graph-traversals:graph}{graph} $G$ with $n \geq 1$ \omterm{def:b1:graph-traversals:graph}{vertices}, the following are equivalent: (i) $G$ is a \omterm{def:b1:graph-traversals:free-tree}{free tree}; (ii) any
two \omterm{def:b1:graph-traversals:graph}{vertices} are joined by exactly one \omterm{def:b1:graph-traversals:path}{simple path}; (iii) $G$ is \omterm{def:b1:graph-traversals:connected}{connected} and has $n - 1$ \omterm{def:b1:graph-traversals:graph}{edges}.
\end{proposition}

\begin{proof}
(i)$\Rightarrow$(ii): a \omterm{def:b1:graph-traversals:path}{simple path} exists, since $G$ is \omterm{def:b1:graph-traversals:connected}{connected} and a shortest \omterm{def:b1:graph-traversals:path}{path} is \omterm{def:b1:graph-traversals:path}{simple}. If two different \omterm{def:b1:graph-traversals:path}{simple
paths} join $u$ and $v$, follow them from $u$ to the first \omterm{def:b1:graph-traversals:graph}{vertex} where they part, $x$, and on the first one to the
first \omterm{def:b1:graph-traversals:graph}{vertex} after $x$ that the second one also visits, $y$; the two pieces from $x$ to $y$ share only their ends and
make a \omterm{def:b1:graph-traversals:path}{cycle}. (ii)$\Rightarrow$(i): $G$ is \omterm{def:b1:graph-traversals:connected}{connected}, and a \omterm{def:b1:graph-traversals:path}{cycle} through $u v$ would give two \omterm{def:b1:graph-traversals:path}{simple paths} from $u$
to $v$, the edge and the rest of the \omterm{def:b1:graph-traversals:path}{cycle}. (i)$\Leftrightarrow$(iii) is \cref{exo:b1:graph-traversals:9}.
\end{proof}

\begin{definition}[Bipartite graph]\label{def:b1:graph-traversals:bipartite}
An undirected \omterm{def:b1:graph-traversals:graph}{graph} is \emph{bipartite}\index{bipartite graph} if its \omterm{def:b1:graph-traversals:graph}{vertices} can be coloured with two colours so that
every edge joins two \omterm{def:b1:graph-traversals:graph}{vertices} of different colours.
\end{definition}

\begin{definition}[Weighted graph, labelled graph]\label{def:b1:graph-traversals:weighted}
A \emph{weighted graph}\index{weighted graph} has a number, its weight, on each edge; a \emph{labelled
graph}\index{labelled graph} has a value of any kind on each edge or \omterm{def:b1:graph-traversals:graph}{vertex}: a road's length, a street's name.
\end{definition}

\section{Representations and their costs}\label{sec:b1:graph-traversals:representations}

The mathematics book stores a \omterm{def:b1:graph-traversals:graph}{graph} in its adjacency matrix: an $n \times n$ table whose entry $(u, v)$ is 1 when $u v$
is an edge. Programs more often keep, for each \omterm{def:b1:graph-traversals:graph}{vertex}, the list of its neighbours.

\begin{definition}[Adjacency list, edge list]\label{def:b1:graph-traversals:lists}
The \emph{adjacency lists}\index{adjacency list} of a \omterm{def:b1:graph-traversals:graph}{graph} give, for each \omterm{def:b1:graph-traversals:graph}{vertex}, the list of its neighbours (in a
\omterm{def:b1:graph-traversals:graph}{directed graph}, the heads of its \omterm{def:b1:graph-traversals:graph}{arcs}). Its \emph{edge list}\index{edge list} is the list of its \omterm{def:b1:graph-traversals:graph}{edges}, as pairs of
\omterm{def:b1:graph-traversals:graph}{vertices}.
\end{definition}

\omcodes{code/bachelor-1/24-graph-traversals/listings/graph.py}{10-32,48-58}{A \omterm{def:b1:graph-traversals:graph}{graph} as \omterm{def:b1:graph-traversals:lists}{adjacency lists}, and the
town's roads loaded}\label{lst:b1:graph-traversals:graph}

\omoutput{code/bachelor-1/24-graph-traversals/outputs/graph.txt}{The two \omterm{def:b1:graph-traversals:graph}{graphs} of the town: each two-way road is two
\omterm{def:b1:graph-traversals:graph}{arcs}, so the roads give $2 \times 53 - 1 = 105$ \omterm{def:b1:graph-traversals:graph}{arcs}.}\label{out:b1:graph-traversals:graph}

The \omterm{def:b1:graph-traversals:degree}{degrees} of the friendship \omterm{def:b1:graph-traversals:graph}{graph} sum to 100, twice its 50 \omterm{def:b1:graph-traversals:graph}{edges}, as the handshake lemma says; the \omterm{def:b1:graph-traversals:degree}{out-degrees} of the
road \omterm{def:b1:graph-traversals:graph}{graph} sum to its number of \omterm{def:b1:graph-traversals:graph}{arcs}.

\begin{center}
\small
\begin{tabular}{l|c|c|c}
 & adjacency matrix & \omterm{def:b1:graph-traversals:lists}{adjacency lists} & \omterm{def:b1:graph-traversals:lists}{edge list} \\ \hline
space & $\Theta(n^2)$ & $\Theta(n + m)$ & $\Theta(m)$ \\
is $u v$ an edge? & $O(1)$ & $O(\deg u)$ & $O(m)$ \\
the neighbours of $u$ & $\Theta(n)$ & $\Theta(\deg u)$ & $\Theta(m)$ \\
add an edge & $O(1)$ & $O(1)$ & $O(1)$ \\
\end{tabular}
\end{center}

\omoutput{code/bachelor-1/24-graph-traversals/outputs/representations.txt}{Numbers each representation stores: the
lists hold one entry per \omterm{def:b1:graph-traversals:graph}{vertex} and one per \omterm{def:b1:graph-traversals:graph}{arc}, twice as many per edge in an undirected
\omterm{def:b1:graph-traversals:graph}{graph}.}\label{out:b1:graph-traversals:representations}

A \omterm{def:b1:graph-traversals:graph}{graph} is \emph{sparse} when $m$ is much smaller than $n^2$, as for roads and friendships, whose \omterm{def:b1:graph-traversals:degree}{degrees} stay small
however large the town; \omterm{def:b1:graph-traversals:lists}{adjacency lists} are then the right choice, and every algorithm of this chapter runs in time
$O(n + m)$ on them. To store a \omterm{def:b1:graph-traversals:graph}{graph} in a file, an \omterm{def:b1:graph-traversals:lists}{edge list} is the simplest: one line per edge.

\omcode{code/bachelor-1/24-graph-traversals/listings/edge_file.py}{9}{26}{A \omterm{def:b1:graph-traversals:graph}{graph} written as an \omterm{def:b1:graph-traversals:lists}{edge list} and read
back}\label{lst:b1:graph-traversals:edge-file}

\omoutput{code/bachelor-1/24-graph-traversals/outputs/edge_file.txt}{The road \omterm{def:b1:graph-traversals:graph}{graph} survives the round trip through a
file.}\label{out:b1:graph-traversals:edge-file}

The first line also lists the \omterm{def:b1:graph-traversals:graph}{vertices}, so that a \omterm{def:b1:graph-traversals:graph}{vertex} without \omterm{def:b1:graph-traversals:graph}{edges}, such as person 40 of the friendship \omterm{def:b1:graph-traversals:graph}{graph},
is not lost.

\section{Depth-first search}\label{sec:b1:graph-traversals:dfs}

\begin{definition}[Depth-first search, discovery and finishing times]\label{def:b1:graph-traversals:dfs}
A \emph{depth-first search}\index{depth-first search} visits a \omterm{def:b1:graph-traversals:graph}{vertex} $u$ by marking it, then visiting in turn each
unmarked neighbour of $u$; on the whole \omterm{def:b1:graph-traversals:graph}{graph}, it visits each \omterm{def:b1:graph-traversals:graph}{vertex} not yet marked. A clock ticks at each start and
each end of a visit: $u$'s \emph{discovery time}\index{discovery time} $d(u)$ is the tick when its visit starts, its
\emph{finishing time}\index{finishing time} $f(u)$ the tick when it ends. Between the two, $u$ is \emph{open}. If the
visit of $u$ calls the visit of $v$, $u$ is the \emph{parent} of $v$; the parent links form the \emph{DFS \omterm{def:b1:graph-traversals:free-tree}{forest}}.
\end{definition}

\begin{definition}[Classes of arcs]\label{def:b1:graph-traversals:edge-classes}
In a \omterm{def:b1:graph-traversals:dfs}{depth-first search} of a \omterm{def:b1:graph-traversals:graph}{directed graph}, an \omterm{def:b1:graph-traversals:graph}{arc} $u \to v$ is a \emph{tree edge}\index{tree edge} if $u$ is the
parent of $v$; a \emph{back edge}\index{back edge} if $v$ is an ancestor of $u$ in the \omterm{def:b1:graph-traversals:free-tree}{forest}, or $u$ itself; a
\emph{forward edge}\index{forward edge} if $v$ is a descendant of $u$ but not its child; a \emph{cross
edge}\index{cross edge} otherwise.
\end{definition}

\omcode{code/bachelor-1/24-graph-traversals/listings/dfs.py}{7}{35}{\omterm{def:b1:graph-traversals:dfs}{Depth-first search}: times, and the class of each
\omterm{def:b1:graph-traversals:graph}{arc}}\label{lst:b1:graph-traversals:dfs}

When the \omterm{def:b1:graph-traversals:graph}{arc} $u \to v$ is examined, $v$ is unmarked (a \omterm{def:b1:graph-traversals:edge-classes}{tree edge}), open (on the chain of visits in progress, so an
ancestor: a \omterm{def:b1:graph-traversals:edge-classes}{back edge}), or finished, and then a descendant of $u$ exactly when it was discovered after $u$.

\omoutput{code/bachelor-1/24-graph-traversals/outputs/dfs.txt}{A small \omterm{def:b1:graph-traversals:graph}{directed graph}: times, classes, and the order
of the iterative search.}\label{out:b1:graph-traversals:dfs}

\begin{omfigure}
\begin{tikzpicture}[v/.style={circle, draw, minimum size=6mm, inner sep=0pt, font=\small},
  t/.style={font=\scriptsize, omDef}, >=stealth]
\node[v] (a) at (0,0) {a}; \node[v] (b) at (2.2,0) {b}; \node[v] (c) at (0,-1.8) {c}; \node[v] (d) at (2.2,-1.8) {d};
\node[v] (e) at (4.6,-1.8) {e}; \node[v] (f) at (4.6,0) {f}; \node[v] (g) at (6.8,-0.9) {g}; \node[v] (h) at (8.8,-0.9)
{h};
\foreach \n/\tm/\pos in {a/1|8/above, b/2|5/above, c/6|7/left, d/3|4/right, e/9|16/below, f/10|15/above,
  g/11|14/above, h/12|13/above} \node[t, \pos=3.3mm] at (\n) {\tm};
\draw[->, very thick] (a) -- (b); \draw[->, very thick] (b) -- (d); \draw[->, very thick] (a) -- (c);
\draw[->, very thick] (e) -- (f); \draw[->, very thick] (f) -- (g); \draw[->, very thick] (g) -- (h);
\draw[->, dashed, omThm, thick] (d) to[bend left=18] (a); \draw[->, dashed, omThm, thick] (g) -- (e);
\draw[->, dotted, thick] (a) to[bend left=18] (d);
\draw[->, gray, thick] (c) -- (b); \draw[->, gray, thick] (e) to[bend left=35] (c);
\node[anchor=west, font=\scriptsize, align=left] at (9.5,-0.3) {thick: tree\\ \textcolor{omThm}{dashed: back}\\
dotted: forward\\ \textcolor{gray}{grey: cross}};
\end{tikzpicture}
\par\smallskip{\small The search of \cref{out:b1:graph-traversals:dfs}, started at \code{a}, then at \code{e}; each
\omterm{def:b1:graph-traversals:graph}{vertex} carries $d|f$. The \omterm{def:b1:graph-traversals:edge-classes}{tree edges} form two trees, rooted at \code{a} and \code{e}.}
\end{omfigure}

\begin{theorem}[Parenthesis theorem]\label{thm:b1:graph-traversals:parenthesis}
For any two \omterm{def:b1:graph-traversals:graph}{vertices} $u$ and $v$ of a \omterm{def:b1:graph-traversals:dfs}{depth-first search}, the intervals $[d(u), f(u)]$ and $[d(v), f(v)]$ are disjoint
or one contains the other; and $[d(v), f(v)] \subset [d(u), f(u)]$ exactly when $v$ is a descendant of $u$, $v \neq
u$.
\end{theorem}

\begin{proof}
Say $d(u) < d(v)$. If $v$ is discovered while $u$ is open, $v$'s visit is called, directly or not, from $u$'s, so $v$
is a descendant of $u$; and a call returns before its caller does, so $f(v) < f(u)$: nested. Otherwise $u$ finished
before $d(v)$: disjoint. Conversely, a descendant of $u$ is discovered during $u$'s visit, so its interval is
nested in $u$'s.
\end{proof}

Writing an opening bracket at each discovery and a closing one at each finish gives a well-bracketed word, hence the
name. A second fact explains what the search finds: when $v$ is discovered, every \omterm{def:b1:graph-traversals:graph}{vertex} reachable from $v$ by a \omterm{def:b1:graph-traversals:path}{path}
of unmarked \omterm{def:b1:graph-traversals:graph}{vertices} becomes a descendant of $v$ (the \emph{\omterm{def:b1:graph-traversals:path}{white-path}} property; take the first \omterm{def:b1:graph-traversals:graph}{vertex} of such a \omterm{def:b1:graph-traversals:path}{path}
that would not be a descendant: its predecessor on the \omterm{def:b1:graph-traversals:path}{path} is one, and examines it before finishing).

\begin{proposition}[Cycles and back edges]\label{prop:b1:graph-traversals:back}
A \omterm{def:b1:graph-traversals:graph}{directed graph} is acyclic if and only if a \omterm{def:b1:graph-traversals:dfs}{depth-first search} finds no \omterm{def:b1:graph-traversals:edge-classes}{back edge}.
\end{proposition}

\begin{proof}
A \omterm{def:b1:graph-traversals:edge-classes}{back edge} $u \to v$, with $v$ an ancestor of $u$, closes a \omterm{def:b1:graph-traversals:path}{cycle} with the tree \omterm{def:b1:graph-traversals:path}{path} from $v$ to $u$. Conversely, let
$C$ be a \omterm{def:b1:graph-traversals:path}{cycle} and $v$ its first \omterm{def:b1:graph-traversals:graph}{vertex} to be discovered, and $u \to v$ the \omterm{def:b1:graph-traversals:graph}{arc} of $C$ that enters $v$. At $d(v)$ the
rest of $C$ is unmarked, so $u$ becomes a descendant of $v$ by the \omterm{def:b1:graph-traversals:path}{white-path} property, and $u \to v$ is a \omterm{def:b1:graph-traversals:edge-classes}{back edge}.
\end{proof}

The recursive search can run out of Python's \omterm{def:b1:c-functions-and-stack:call-stack}{call stack} on a long \omterm{def:b1:graph-traversals:path}{path}. An explicit stack does the same work in a
loop; pushing all the neighbours at once, it discovers the \omterm{def:b1:graph-traversals:graph}{vertices} in the same order, but it no longer knows when a
\omterm{def:b1:graph-traversals:graph}{vertex} finishes. For \omterm{def:b1:graph-traversals:dfs}{finishing times}, the stack holds each open \omterm{def:b1:graph-traversals:graph}{vertex} with the \omterm{def:b1:linked-structures:iterator}{iterator} over its remaining \omterm{def:b1:graph-traversals:graph}{arcs}
(\cref{lst:b1:graph-traversals:kosaraju}).

\omcode{code/bachelor-1/24-graph-traversals/listings/dfs.py}{38}{48}{The order of discovery, with a
stack}\label{lst:b1:graph-traversals:dfs-iterative}

\section{Breadth-first search}\label{sec:b1:graph-traversals:bfs}

\begin{definition}[Breadth-first search, BFS tree]\label{def:b1:graph-traversals:bfs}
A \emph{breadth-first search}\index{breadth-first search} from $s$ keeps a queue of discovered \omterm{def:b1:graph-traversals:graph}{vertices}, initially $s$;
it removes the \omterm{def:b1:graph-traversals:graph}{vertex} $u$ at the head and discovers each undiscovered neighbour $v$ of $u$, setting $\code{dist}[v] =
\code{dist}[u] + 1$ and the parent of $v$ to $u$, and adds it at the tail. The parent links form the \emph{BFS
tree}\index{BFS tree}, rooted at $s$.
\end{definition}

\omcode{code/bachelor-1/24-graph-traversals/listings/bfs.py}{8}{19}{\omterm{def:b1:graph-traversals:bfs}{Breadth-first search} with a
queue}\label{lst:b1:graph-traversals:bfs}

\begin{theorem}[Breadth-first search computes distances]\label{thm:b1:graph-traversals:bfs}
At the end of a \omterm{def:b1:graph-traversals:bfs}{breadth-first search} from $s$, the discovered \omterm{def:b1:graph-traversals:graph}{vertices} are those reachable from $s$, $\code{dist}[v] =
\delta(s, v)$ for each, and the \omterm{def:b1:graph-traversals:path}{path} of parents from $v$ up to $s$, read backwards, is a shortest \omterm{def:b1:graph-traversals:path}{path}.
\end{theorem}

\begin{proof}
Invariant: the values of \code{dist} in the queue, from head to tail, are non-decreasing and differ by at most 1. It
holds for $[s]$; removing the head $u$, of value $k$, and appending \omterm{def:b1:graph-traversals:graph}{vertices} of value $k + 1$ keeps it, since the other
values were between $k$ and $k + 1$. So \omterm{def:b1:graph-traversals:graph}{vertices} leave the queue by non-decreasing \code{dist}. Each $\code{dist}[v]$
is the length of the parent \omterm{def:b1:graph-traversals:path}{path}, so $\code{dist}[v] \geq \delta(s, v)$. Suppose some $v$ has $\code{dist}[v] >
\delta(s, v) = k$, with $k$ smallest; $k \geq 1$, and let $u$ be the \omterm{def:b1:graph-traversals:graph}{vertex} before $v$ on a shortest \omterm{def:b1:graph-traversals:path}{path}, so
$\code{dist}[u] = \delta(s, u) = k - 1$. When $u$ left the queue, $v$ was discovered then, with value $k$, or before,
from a \omterm{def:b1:graph-traversals:graph}{vertex} that left the queue no later than $u$, with value at most $k$: a contradiction. Reachable \omterm{def:b1:graph-traversals:graph}{vertices} are
discovered by the same argument, with $\infty$ in place of a value.
\end{proof}

\omoutput{code/bachelor-1/24-graph-traversals/outputs/bfs.txt}{The places by number of roads from the Town Square, and a
farthest one.}\label{out:b1:graph-traversals:bfs}

\begin{omfigure}
\begin{tikzpicture}
\begin{axis}[omrecord, width=10cm, height=8.6cm, xlabel={$x$ (km)}, ylabel={$y$ (km)}, xmin=-0.2, xmax=6, ymin=-0.2,
    ymax=5.1, axis equal image, grid=none]
  \addplot[black!25, line width=0.6pt] table[x=x, y=y]{code/bachelor-1/24-graph-traversals/figures/roads.dat};
  \addplot[omThm, line width=1.2pt, -{Stealth[length=2.5mm]}] table[x=x, y=y]
    {code/bachelor-1/24-graph-traversals/figures/oneway.dat};
  \addplot[only marks, mark=*, mark size=1.6pt, omDef, nodes near coords, point meta=explicit symbolic,
    nodes near coords style={font=\tiny, black!70, anchor=south west, inner sep=1pt}]
    table[x=x, y=y, meta=d]{code/bachelor-1/24-graph-traversals/figures/places.dat};
  \node[font=\scriptsize, anchor=north east] at (axis cs:4.27,1.40) {Town Square};
\end{axis}
\end{tikzpicture}
\par\smallskip{\small Brenmoor's roads, drawn from the coordinates of the places by \code{figures/town_graphs.py}; the
one-way road is the arrow. Each place carries its \omterm{def:b1:graph-traversals:distance}{distance} in roads from the Town Square: the layers of
\cref{out:b1:graph-traversals:bfs}.}
\end{omfigure}

\omterm{def:b1:graph-traversals:bfs}{Breadth-first search} counts roads, not metres: the route it finds to the Hospital has the fewest roads, not the fewest
kilometres. Lengths are the business of \cref{ch:b1:shortest-paths}.

The same search two-colours a \omterm{def:b1:graph-traversals:graph}{graph}: colour $s$ with 0, and each newly discovered \omterm{def:b1:graph-traversals:graph}{vertex} with the colour its parent
does not have. If an edge joins two \omterm{def:b1:graph-traversals:graph}{vertices} of the same colour, the \omterm{def:b1:graph-traversals:graph}{graph} is not \omterm{def:b1:graph-traversals:bipartite}{bipartite}, and the edge with the two
tree \omterm{def:b1:graph-traversals:path}{paths} up to the common ancestor of its ends is an odd \omterm{def:b1:graph-traversals:path}{cycle}: the two \omterm{def:b1:graph-traversals:path}{paths} have lengths of the same parity, since
the colours of their ends agree.

\omcode{code/bachelor-1/24-graph-traversals/listings/bipartite.py}{8}{36}{Two colours, or an odd
\omterm{def:b1:graph-traversals:path}{cycle}}\label{lst:b1:graph-traversals:bipartite}

\omoutput{code/bachelor-1/24-graph-traversals/outputs/bipartite.txt}{The friendship \omterm{def:b1:graph-traversals:graph}{graph} is not \omterm{def:b1:graph-traversals:bipartite}{bipartite}: three
people are friends with one another.}\label{out:b1:graph-traversals:bipartite}

An odd \omterm{def:b1:graph-traversals:path}{cycle} is a \emph{certificate}: anyone can check in a moment that it is a \omterm{def:b1:graph-traversals:path}{cycle} of the \omterm{def:b1:graph-traversals:graph}{graph} and that its length
is odd, and no \omterm{def:b1:propositional-logic:colouring}{two-colouring} survives it (\cref{exo:b1:graph-traversals:12}).

\section{Applications}\label{sec:b1:graph-traversals:applications}

\paragraph{\omterm{def:b1:graph-traversals:connected}{Connected components}.} One search from an unmarked \omterm{def:b1:graph-traversals:graph}{vertex} marks exactly its component; restarting from
each unmarked \omterm{def:b1:graph-traversals:graph}{vertex} in turn finds them all, in time $O(n + m)$.

\omcode{code/bachelor-1/24-graph-traversals/listings/components.py}{6}{22}{The \omterm{def:b1:graph-traversals:connected}{connected
components}}\label{lst:b1:graph-traversals:components}

\omoutput{code/bachelor-1/24-graph-traversals/outputs/components.txt}{The friendship \omterm{def:b1:graph-traversals:graph}{graph} has three components; the
third is person 40 alone.}\label{out:b1:graph-traversals:components}

\begin{omfigure}
\begin{tikzpicture}
\begin{axis}[hide axis, width=10.5cm, height=9.5cm, axis equal image, xmin=-2.8, xmax=9.8, ymin=-8.6, ymax=2.8]
  \addplot[black!35, line width=0.5pt] table[x=x, y=y]{code/bachelor-1/24-graph-traversals/figures/friend_edges.dat};
  \addplot[omThm, line width=1.3pt] table[x=x, y=y]{code/bachelor-1/24-graph-traversals/figures/friend_bridges.dat};
  \addplot[only marks, mark=*, mark size=1.4pt, omDef, nodes near coords, point meta=explicit symbolic,
    nodes near coords style={font=\tiny, black!80, anchor=south west, inner sep=0.5pt, fill=white}]
    table[x=x, y=y, meta=id]{code/bachelor-1/24-graph-traversals/figures/people.dat};
\end{axis}
\end{tikzpicture}
\par\smallskip{\small The friendship \omterm{def:b1:graph-traversals:graph}{graph}, drawn by \code{figures/town_graphs.py}: people 1--14, 15--27 and 28--39 on
three circles, person 40 apart; the \omterm{def:b1:graph-traversals:bridge}{bridges} are thick.}
\end{omfigure}

\paragraph{\omterm{def:b1:graph-traversals:topo}{Topological order}.}

\begin{definition}[Topological order]\label{def:b1:graph-traversals:topo}
A \emph{topological order}\index{topological order} of a \omterm{def:b1:graph-traversals:graph}{directed graph} is an order of its \omterm{def:b1:graph-traversals:graph}{vertices} in which, for every
\omterm{def:b1:graph-traversals:graph}{arc} $u \to v$, $u$ comes before $v$.
\end{definition}

A \omterm{def:b1:graph-traversals:graph}{graph} with a \omterm{def:b1:graph-traversals:path}{cycle} has none: the first \omterm{def:b1:graph-traversals:graph}{vertex} of the \omterm{def:b1:graph-traversals:path}{cycle} in the order would come after its predecessor on the
\omterm{def:b1:graph-traversals:path}{cycle}. A DAG has one, and \omterm{def:b1:graph-traversals:dfs}{depth-first search} finds it.

\begin{proposition}[Order by decreasing finishing time]\label{prop:b1:graph-traversals:topo}
In a \omterm{def:b1:graph-traversals:dfs}{depth-first search} of a DAG, $f(v) < f(u)$ for every \omterm{def:b1:graph-traversals:graph}{arc} $u \to v$. So the \omterm{def:b1:graph-traversals:graph}{vertices} by decreasing \omterm{def:b1:graph-traversals:dfs}{finishing time}
are in \omterm{def:b1:graph-traversals:topo}{topological order}.
\end{proposition}

\begin{proof}
When $u \to v$ is examined, $v$ is not open: it would be an ancestor of $u$ and the \omterm{def:b1:graph-traversals:graph}{arc} a \omterm{def:b1:graph-traversals:edge-classes}{back edge}, which a DAG does
not have (\cref{prop:b1:graph-traversals:back}). If $v$ is unmarked, it becomes a descendant of $u$ and finishes
first (\cref{thm:b1:graph-traversals:parenthesis}); if it is finished, it finished before $u$, which is still open.
\end{proof}

Kahn's method works from the other end: a DAG has a \omterm{def:b1:graph-traversals:graph}{vertex} with no incoming \omterm{def:b1:graph-traversals:graph}{arc} (follow \omterm{def:b1:graph-traversals:graph}{arcs} backwards from any \omterm{def:b1:graph-traversals:graph}{vertex};
without a \omterm{def:b1:graph-traversals:path}{cycle} the walk must stop). Output one, remove it with its \omterm{def:b1:graph-traversals:graph}{arcs}, and repeat. With \omterm{def:b1:graph-traversals:degree}{in-degrees} counted once and
decreased as \omterm{def:b1:graph-traversals:graph}{vertices} go, it runs in $O(n + m)$; if \omterm{def:b1:graph-traversals:graph}{vertices} remain that all have incoming \omterm{def:b1:graph-traversals:graph}{arcs}, there is a \omterm{def:b1:graph-traversals:path}{cycle}.

\omcode{code/bachelor-1/24-graph-traversals/listings/topo.py}{17}{47}{A \omterm{def:b1:graph-traversals:topo}{topological order}, by \omterm{def:b1:graph-traversals:dfs}{depth-first search} and by
Kahn's method}\label{lst:b1:graph-traversals:topo}

\omoutput{code/bachelor-1/24-graph-traversals/outputs/topo.txt}{Some chapters of this book in an order that respects
which relies on which: two different orders, both valid.}\label{out:b1:graph-traversals:topo}

\paragraph{\omterm{def:b1:graph-traversals:connected}{Strongly connected components}.}

\begin{definition}[Component graph]\label{def:b1:graph-traversals:condensation}
The \emph{component graph}\index{component graph}, or \emph{condensation}\index{condensation}, of a \omterm{def:b1:graph-traversals:graph}{directed graph} has
its \omterm{def:b1:graph-traversals:connected}{strongly connected components} as \omterm{def:b1:graph-traversals:graph}{vertices}, and an \omterm{def:b1:graph-traversals:graph}{arc} $C \to C'$ when some \omterm{def:b1:graph-traversals:graph}{arc} of the \omterm{def:b1:graph-traversals:graph}{graph} goes from $C$ to
$C' \neq C$.
\end{definition}

The \omterm{def:b1:graph-traversals:condensation}{component graph} is a DAG: a \omterm{def:b1:graph-traversals:path}{cycle} through $C$ and $C'$ would make their \omterm{def:b1:graph-traversals:graph}{vertices} reachable from one another, so
$C = C'$. A search started in a component with no \omterm{def:b1:graph-traversals:graph}{arc} leaving it in the \omterm{def:b1:graph-traversals:condensation}{component graph} marks exactly that component;
the difficulty is to find such a component first.

\begin{method}[Kosaraju's algorithm]\label{met:b1:graph-traversals:kosaraju}
\begin{enumerate}
\item Run a \omterm{def:b1:graph-traversals:dfs}{depth-first search} of $G$ and list the \omterm{def:b1:graph-traversals:graph}{vertices} by decreasing \omterm{def:b1:graph-traversals:dfs}{finishing time}.
\item Reverse every \omterm{def:b1:graph-traversals:graph}{arc} of $G$.
\item In the reversed \omterm{def:b1:graph-traversals:graph}{graph}, take the \omterm{def:b1:graph-traversals:graph}{vertices} in that list; from each one not yet assigned, search, and assign every
\omterm{def:b1:graph-traversals:graph}{vertex} reached that is not yet assigned to a new component.
\end{enumerate}
\end{method}

\omcode{code/bachelor-1/24-graph-traversals/listings/kosaraju.py}{6}{43}{Kosaraju's algorithm, with an iterative
search for the finishing order}\label{lst:b1:graph-traversals:kosaraju}

\omoutput{code/bachelor-1/24-graph-traversals/outputs/kosaraju.txt}{The small \omterm{def:b1:graph-traversals:graph}{graph} of the figure below;
the roads with their one one-way street, and with every road one-way in the direction the table lists
it.}\label{out:b1:graph-traversals:kosaraju}

The answer to the first question is yes: with road 1 one-way, the thirty places form a single \omterm{def:b1:graph-traversals:connected}{strongly connected
component}. The second question needs care: making each road one-way in the direction the table happens to list it
leaves every place alone in its component, since the table lists each road from the smaller place number to the
larger, which makes a DAG.

\begin{lemma}[Finishing times and components]\label{lem:b1:graph-traversals:finish}
If $C$ and $C'$ are distinct \omterm{def:b1:graph-traversals:connected}{strongly connected components} and some \omterm{def:b1:graph-traversals:graph}{arc} goes from $C$ to $C'$, the largest \omterm{def:b1:graph-traversals:dfs}{finishing
time} in $C$ is larger than the largest in $C'$.
\end{lemma}

\begin{proof}
Let $x$ be the first \omterm{def:b1:graph-traversals:graph}{vertex} of $C \cup C'$ discovered. If $x \in C$, every \omterm{def:b1:graph-traversals:graph}{vertex} of $C \cup C'$ is reachable from $x$
through unmarked \omterm{def:b1:graph-traversals:graph}{vertices}, so all are descendants of $x$ and finish before it. If $x \in C'$, no \omterm{def:b1:graph-traversals:graph}{vertex} of $C$ is
reachable from $x$ (or $C$ and $C'$ would merge), so all of $C'$ finishes, as descendants of $x$, before any \omterm{def:b1:graph-traversals:graph}{vertex} of
$C$ is discovered.
\end{proof}

The proof that Kosaraju's algorithm is correct is \cref{exo:b1:graph-traversals:10}: the \omterm{def:b1:graph-traversals:graph}{vertex} with the largest
\omterm{def:b1:graph-traversals:dfs}{finishing time} lies in a component with no \omterm{def:b1:graph-traversals:graph}{arc} entering it, and in the reversed \omterm{def:b1:graph-traversals:graph}{graph} that component has no \omterm{def:b1:graph-traversals:graph}{arc}
leaving it. Both searches cost $O(n + m)$, and so does reversing the \omterm{def:b1:graph-traversals:graph}{graph}.

\begin{omfigure}
\begin{tikzpicture}[v/.style={circle, draw, minimum size=6mm, inner sep=0pt, font=\small}, >=stealth,
  comp/.style={rounded corners=3mm, fill=omDef!12, draw=omDef!50}]
\fill[comp] (-0.5,-2.1) rectangle (2.7,0.5); \fill[comp] (3.4,-2.1) rectangle (6.0,-1.1);
\fill[comp] (3.4,-0.4) rectangle (6.0,0.5); \fill[comp] (6.7,-2.0) rectangle (7.6,-1.2);
\node[v] (a) at (0,0) {a}; \node[v] (b) at (2.2,0) {b}; \node[v] (c) at (1.1,-1.6) {c};
\node[v] (d) at (3.9,0.05) {d}; \node[v] (e) at (5.5,0.05) {e};
\node[v] (f) at (3.9,-1.6) {f}; \node[v] (g) at (5.5,-1.6) {g}; \node[v] (h) at (7.15,-1.6) {h};
\foreach \u/\w in {a/b, b/c, c/a, b/d, c/f, e/g, g/h} \draw[->, thick] (\u) -- (\w);
\draw[->, thick] (d) to[bend left=20] (e); \draw[->, thick] (e) to[bend left=20] (d);
\draw[->, thick] (f) to[bend left=20] (g); \draw[->, thick] (g) to[bend left=20] (f);
\begin{scope}[xshift=90mm, yshift=-8mm, n/.style={draw, rounded corners=2mm, fill=omDef!12, font=\small,
  minimum height=6mm}]
\node[n] (A) at (0,0) {abc}; \node[n] (D) at (1.6,0.9) {de}; \node[n] (F) at (1.6,-0.9) {fg}; \node[n] (H) at (3.2,-0.9)
{h};
\draw[->, thick] (A) -- (D); \draw[->, thick] (A) -- (F); \draw[->, thick] (D) -- (F); \draw[->, thick] (F) -- (H);
\end{scope}
\end{tikzpicture}
\par\smallskip{\small Left, a \omterm{def:b1:graph-traversals:graph}{directed graph} and its four \omterm{def:b1:graph-traversals:connected}{strongly connected components} (shaded); right, its
\omterm{def:b1:graph-traversals:condensation}{component graph}, a DAG. Kosaraju's second search, on the reversed \omterm{def:b1:graph-traversals:graph}{arcs}, starts in \code{abc}, which no \omterm{def:b1:graph-traversals:graph}{arc} enters:
reversed, no \omterm{def:b1:graph-traversals:graph}{arc} leaves it, so the search stays inside.}
\end{omfigure}

\paragraph{\omterm{def:b1:graph-traversals:bridge}{Bridges}.}

\begin{definition}[Bridge]\label{def:b1:graph-traversals:bridge}
A \emph{bridge}\index{bridge} of an undirected \omterm{def:b1:graph-traversals:graph}{graph} is an edge whose removal increases the number of \omterm{def:b1:graph-traversals:connected}{connected
components}.
\end{definition}

An edge is a \omterm{def:b1:graph-traversals:bridge}{bridge} exactly when it lies on no \omterm{def:b1:graph-traversals:path}{cycle}: if $uv$ is on a \omterm{def:b1:graph-traversals:path}{cycle}, the rest of the \omterm{def:b1:graph-traversals:path}{cycle} still joins $u$
and $v$; if it is not, no other \omterm{def:b1:graph-traversals:path}{path} joins them. In a \omterm{def:b1:graph-traversals:dfs}{depth-first search} of an undirected \omterm{def:b1:graph-traversals:graph}{graph} there are no \omterm{def:b1:graph-traversals:edge-classes}{cross
edges} (an edge to a finished \omterm{def:b1:graph-traversals:graph}{vertex} not a descendant would have been taken from the earlier one), so every non-tree
edge joins a \omterm{def:b1:graph-traversals:graph}{vertex} to an ancestor. Let $\mathrm{low}(v)$ be the smallest \omterm{def:b1:graph-traversals:dfs}{discovery time} reachable from the \omterm{def:b1:induction-and-termination:rooted-tree}{subtree} of
$v$ by going down \omterm{def:b1:graph-traversals:edge-classes}{tree edges} and then along at most one other edge. A \omterm{def:b1:graph-traversals:edge-classes}{tree edge} $u v$, with $u$ the parent, is on a
\omterm{def:b1:graph-traversals:path}{cycle} exactly when some edge leaves the \omterm{def:b1:induction-and-termination:rooted-tree}{subtree} of $v$ for $u$ or above, that is when $\mathrm{low}(v) \leq d(u)$. The
low points are computed on the way back up, in one search.

\begin{theorem}[Robbins' theorem]\label{thm:b1:graph-traversals:robbins}
The \omterm{def:b1:graph-traversals:graph}{edges} of a \omterm{def:b1:graph-traversals:connected}{connected} undirected \omterm{def:b1:graph-traversals:graph}{graph} can be oriented so that the resulting \omterm{def:b1:graph-traversals:graph}{directed graph} is \omterm{def:b1:graph-traversals:connected}{strongly connected}
if and only if the \omterm{def:b1:graph-traversals:graph}{graph} has no \omterm{def:b1:graph-traversals:bridge}{bridge}.
\end{theorem}

Robbins proved it in 1939, in a paper whose title names its application, traffic control. The proof, the \omterm{def:b1:graph-traversals:bridge}{bridges} and
the orientation of Brenmoor's roads are the weekend problem.

\section{Exercises}\label{sec:b1:graph-traversals:exercises}

\begin{exercise}[$\star$]\label{exo:b1:graph-traversals:1}
A \omterm{def:b1:graph-traversals:graph}{graph} on the \omterm{def:b1:graph-traversals:graph}{vertices} 1 to 6 has the \omterm{def:b1:graph-traversals:graph}{edges} $12, 13, 23, 24, 35, 45, 56$. Give the \omterm{def:b1:graph-traversals:degree}{degrees}, check the handshake lemma,
and list the \omterm{def:b1:graph-traversals:graph}{vertices} of odd \omterm{def:b1:graph-traversals:degree}{degree}.
\end{exercise}

\begin{exercise}[$\star$]\label{exo:b1:graph-traversals:2}
Run \code{dfs} on the \omterm{def:b1:graph-traversals:graph}{arcs} $1 \to 2$, $1 \to 3$, $2 \to 4$, $3 \to 4$, $4 \to 1$, $3 \to 5$, $5 \to 6$, $6 \to 5$, added
in this order. Give the times and the class of each \omterm{def:b1:graph-traversals:graph}{arc}.
\end{exercise}

\begin{exercise}[$\star$]\label{exo:b1:graph-traversals:3}
Give the \omterm{def:b1:graph-traversals:distance}{distances} from \omterm{def:b1:graph-traversals:graph}{vertex} 1 in the \omterm{def:b1:graph-traversals:graph}{graph} of \cref{exo:b1:graph-traversals:1}, and a \omterm{def:b1:graph-traversals:bfs}{BFS tree}.
\end{exercise}

\begin{exercise}[$\star$]\label{exo:b1:graph-traversals:4}
Dinner: boiling water comes before cooking pasta, chopping onions before frying them, the fried onions and an opened tin
before the sauce, the pasta and the sauce before serving. Give a \omterm{def:b1:graph-traversals:topo}{topological order} by Kahn's method and one by
\omterm{def:b1:graph-traversals:dfs}{depth-first search}.
\end{exercise}

\begin{exercise}[$\star\star$]\label{exo:b1:graph-traversals:5}
This search marks a \omterm{def:b1:graph-traversals:graph}{vertex} only when its visit ends. Explain the two lines of output; what happens on a \omterm{def:b1:graph-traversals:graph}{graph} with
a \omterm{def:b1:graph-traversals:path}{cycle}, and why does a DAG survive?
\omcode{code/bachelor-1/24-graph-traversals/listings/exo_mark_exit.py}{8}{14}{Marked on exit}
\omoutput[2-3]{code/bachelor-1/24-graph-traversals/outputs/exo_mark_exit.txt}{A chain of ten diamonds, and a
\omterm{def:b1:graph-traversals:graph}{two-vertex} \omterm{def:b1:graph-traversals:path}{cycle}.}
\end{exercise}

\begin{exercise}[$\star\star$]\label{exo:b1:graph-traversals:6}
Find the \omterm{def:b1:graph-traversals:graph}{graph} this test misjudges, and repair it.
\omcode{code/bachelor-1/24-graph-traversals/listings/exo_bipartite_one.py}{8}{19}{A \omterm{def:b1:graph-traversals:bipartite}{bipartite} test}
\omoutput[2-2]{code/bachelor-1/24-graph-traversals/outputs/exo_bipartite_one.txt}{A \omterm{def:b1:graph-traversals:path}{path} 1--2--3--4 and a triangle 5,
6, 7.}
\end{exercise}

\begin{exercise}[$\star\star$]\label{exo:b1:graph-traversals:7}
Modify \omterm{def:b1:graph-traversals:bfs}{breadth-first search} to count, for every \omterm{def:b1:graph-traversals:graph}{vertex} $v$, the shortest \omterm{def:b1:graph-traversals:path}{paths} from $s$ to $v$, in time $O(n + m)$.
\end{exercise}

\begin{exercise}[$\star\star$]\label{exo:b1:graph-traversals:8}
Write a function that returns a \omterm{def:b1:graph-traversals:path}{cycle} of a \omterm{def:b1:graph-traversals:graph}{directed graph}, as a list of \omterm{def:b1:graph-traversals:graph}{vertices}, or \code{None} if the \omterm{def:b1:graph-traversals:graph}{graph} is
acyclic, in time $O(n + m)$.
\end{exercise}

\begin{exercise}[$\star\star$]\label{exo:b1:graph-traversals:9}
Show that a \omterm{def:b1:graph-traversals:free-tree}{free tree} with $n$ \omterm{def:b1:graph-traversals:graph}{vertices} has $n - 1$ \omterm{def:b1:graph-traversals:graph}{edges}, and that a \omterm{def:b1:graph-traversals:connected}{connected graph} with $n$ \omterm{def:b1:graph-traversals:graph}{vertices} and $n - 1$
\omterm{def:b1:graph-traversals:graph}{edges} is a \omterm{def:b1:graph-traversals:free-tree}{free tree}.
\end{exercise}

\begin{exercise}[$\star\star\star$]\label{exo:b1:graph-traversals:10}
Prove that Kosaraju's algorithm finds the \omterm{def:b1:graph-traversals:connected}{strongly connected components}.
\end{exercise}

\begin{exercise}[$\star\star\star$]\label{exo:b1:graph-traversals:11}
An \emph{Euler circuit} goes along every edge of a \omterm{def:b1:graph-traversals:graph}{graph} exactly once and returns to its start. Write a function that
finds one in a \omterm{def:b1:graph-traversals:connected}{connected graph} whose \omterm{def:b1:graph-traversals:degree}{degrees} are all even, in time $O(n + m)$, and test it.
\end{exercise}

\begin{exercise}[$\star\star\star$]\label{exo:b1:graph-traversals:12}
Prove that an undirected \omterm{def:b1:graph-traversals:graph}{graph} is \omterm{def:b1:graph-traversals:bipartite}{bipartite} if and only if it has no \omterm{def:b1:graph-traversals:path}{cycle} of odd length.
\end{exercise}

\section{Problem: Robbins' town}\label{sec:b1:graph-traversals:problem}

\begin{problem}[{Weekend problem --- Robbins' town}]\label{pb:b1:graph-traversals:1}
The town council of Brenmoor wants every road one-way, and every place still reachable from every other. Before the
roads, the friendship \omterm{def:b1:graph-traversals:graph}{graph} serves as a warm-up.
\omcode{code/bachelor-1/24-graph-traversals/listings/bridges.py}{7}{32}{\omterm{def:b1:graph-traversals:bridge}{Bridges} by low
points}\label{lst:b1:graph-traversals:bridges}
\omoutput{code/bachelor-1/24-graph-traversals/outputs/bridges.txt}{\omterm{def:b1:graph-traversals:bridge}{Bridges} among friends, and among
roads.}\label{out:b1:graph-traversals:bridges}
\omcode{code/bachelor-1/24-graph-traversals/listings/orient.py}{9}{25}{Orienting the roads by a \omterm{def:b1:graph-traversals:dfs}{depth-first
search}}\label{lst:b1:graph-traversals:orient}
\omoutput{code/bachelor-1/24-graph-traversals/outputs/orient.txt}{Every road one-way, and the town still \omterm{def:b1:graph-traversals:connected}{strongly
connected}.}\label{out:b1:graph-traversals:orient}

\needspace{4\baselineskip}
\medskip\noindent\textbf{Part I --- The friendship graph.}
\begin{enumerate}
\item Using \cref{out:b1:graph-traversals:graph}, how many friendships does person 40 have, and what is the average
\omterm{def:b1:graph-traversals:degree}{degree}?
\item Why does \cref{out:b1:graph-traversals:components} count person 40 as a component?
\item What do the components tell about the three groups of friends drawn in \cref{sec:b1:graph-traversals:applications}?
\item Why does the triangle of \cref{out:b1:graph-traversals:bipartite} prove that the \omterm{def:b1:graph-traversals:graph}{graph} is not \omterm{def:b1:graph-traversals:bipartite}{bipartite}?
\item Could two of the components be \omterm{def:b1:graph-traversals:bipartite}{bipartite} while the whole \omterm{def:b1:graph-traversals:graph}{graph} is not? What would the program print?
\item What is the cost of the component and \omterm{def:b1:graph-traversals:bipartite}{bipartite} searches, in $n$ and $m$?
\end{enumerate}

\needspace{4\baselineskip}
\medskip\noindent\textbf{Part II --- The roads and the one-way street.}
\begin{enumerate}[resume]
\item How many \omterm{def:b1:graph-traversals:graph}{arcs} would the road \omterm{def:b1:graph-traversals:graph}{graph} have if road 1 were two-way?
\item Explain the single component of \cref{out:b1:graph-traversals:kosaraju}: which \omterm{def:b1:graph-traversals:path}{path} leads from the Library back
to the Town Square?
\item Why does listing each road from its smaller place number make a DAG, and why does a DAG have only components of one
\omterm{def:b1:graph-traversals:graph}{vertex}?
\item In the \omterm{def:b1:graph-traversals:condensation}{component graph} of the roads oriented as listed, which \omterm{def:b1:graph-traversals:graph}{vertices} have no incoming \omterm{def:b1:graph-traversals:graph}{arc}?
\item What is the cost of Kosaraju's algorithm on the roads?
\item Why does reversing every \omterm{def:b1:graph-traversals:graph}{arc} leave the \omterm{def:b1:graph-traversals:connected}{strongly connected components} unchanged?
\end{enumerate}

\needspace{4\baselineskip}
\medskip\noindent\textbf{Part III --- Bridges.}
\begin{enumerate}[resume]
\item Show that an edge is a \omterm{def:b1:graph-traversals:bridge}{bridge} if and only if it lies on no \omterm{def:b1:graph-traversals:path}{cycle}.
\item Why does a \omterm{def:b1:graph-traversals:dfs}{depth-first search} of an undirected \omterm{def:b1:graph-traversals:graph}{graph} have no \omterm{def:b1:graph-traversals:edge-classes}{cross edges}?
\item Prove that a \omterm{def:b1:graph-traversals:edge-classes}{tree edge} $u v$, $u$ the parent, is a \omterm{def:b1:graph-traversals:bridge}{bridge} if and only if $\mathrm{low}(v) > d(u)$.
\item Why does \code{visit} skip the edge back to the parent, and only once?
\item The friendship \omterm{def:b1:graph-traversals:graph}{graph} has 12 \omterm{def:b1:graph-traversals:bridge}{bridges}: what do they say about the groups of friends?
\item Why can the road \omterm{def:b1:graph-traversals:graph}{graph} have no \omterm{def:b1:graph-traversals:bridge}{bridge}, given how its roads were laid (each place joined to its three nearest)? Is
this a proof?
\end{enumerate}

\needspace{4\baselineskip}
\medskip\noindent\textbf{Part IV --- Robbins' theorem.}
\begin{enumerate}[resume]
\item Show that a \omterm{def:b1:graph-traversals:graph}{graph} with a \omterm{def:b1:graph-traversals:bridge}{bridge} has no \omterm{def:b1:graph-traversals:connected}{strongly connected} orientation.
\item In \code{orient}, why is every edge oriented exactly once?
\item Show that in the orientation, every \omterm{def:b1:graph-traversals:graph}{vertex} is reachable from the root.
\item Show that, if the \omterm{def:b1:graph-traversals:graph}{graph} has no \omterm{def:b1:graph-traversals:bridge}{bridge}, the root is reachable from every \omterm{def:b1:graph-traversals:graph}{vertex} in the orientation.
\item Deduce Robbins' theorem.
\item Why does road 1 keep the direction it has today?
\item State the named result, \emph{Robbins' theorem}, and what the program checked on Brenmoor.
\end{enumerate}
\end{problem}
